% 原书分章 PDF 第 6—7 页：第 28.3 节。
\section{超级学习器}

TMLE 需要似然各因子的初始估计。在特定处理下均值的例子中，我们需要估计 $\bar Q(A,W)$ 与 $\bar g(W)$。在靶向学习框架内，通常使用超级学习器估计这些因子，第 3 章及其他部分已经讨论过这一方法。本章采用离散超级学习器（discrete super learner），也称交叉验证选择器，它选择使风险最小的单个算法。

基于损失的估计使我们能够客观评价估计的质量，并据此在相互竞争的估计量中进行选择。超级学习器是这一框架的一种具体实现。要理解超级学习器在非参数自助样本上的行为，理解这一框架十分重要。在超级学习器的情境中，我们把似然各因子的估计量称为“学习器”。下面着重讨论学习 $\bar Q_0(P_0)=E_{P_0}[Y\mid A,W]$ 的估计这一例子，但同样的讨论也适用于似然的其他因子。这里，$\bar Q(P)$ 表示基于 $P$ 得到的 $\bar Q$ 的估计。考虑从可能分布的类 $\overline{\mathbf Q}$ 中选择一个估计 $\bar Q$ 的问题。在离散超级学习器中，这变为从候选学习器集合
$\bar Q_k:k=1,\ldots,K$ 中进行选择。

这一框架的关键是损失函数 $L(\bar Q,O)$，它描述学习器的预测值与真实观测值之间差异的严重程度。例如，平方误差损失为
$L(\bar Q,O)=(\bar Q-Y)^2$。由此得到风险，即损失关于分布 $P$ 的期望：
\[
 \theta(\bar Q,P)=PL(\bar Q,O)=E_{P_0}[L(\bar Q,O)].
\]
在真实分布 $P_0$ 处评价，得到真实风险
$\theta_0(\bar Q)=\theta(\bar Q,P_0)$。它提供了选择学习器的准则：
\[
 \bar Q_0=\argmin_{\bar Q\in\overline{\mathbf Q}}
 \theta_0(\bar Q)=\bar Q(P_0),
\]
即我们所需的学习器是使真实风险最小的学习器。关键在于，它应等于我们希望估计的参数，此处即 $\bar Q_0(A,W)$。风险在极小化解处的值称为最优风险（optimal risk），或不可约误差（irreducible error）：
\[
 \theta_0(\bar Q_0)=\min_{\bar Q\in\overline{\mathbf Q}}\theta_0(\bar Q).
\]
于是，在离散超级学习器中，可将预言选择器定义为
\[
 \tilde k_n=\argmin_{k\in\{1,\ldots,K_n\}}\theta_0(\bar Q_k),
\]
它从候选集合中选择使真实风险最小的学习器。如果已知 $P_0$，这就是我们会选择的学习器。

经验风险估计或回代风险估计（resubstitution risk estimate）
\[
 \hat\theta_{P_n}(\bar Q(P_n))=\theta(\bar Q(P_n),P_n)
\]
在用于训练学习器的同一数据集上估计风险。众所周知，在多数情况下，它过于乐观，即存在向下偏差，并且这种乐观程度随模型复杂度增加。因此，回代选择器
$\argmin_{k\in\{1,\ldots,K_n\}}\hat\theta_{P_n}(\bar Q_k(P_n))$
会选择对数据“过拟合”的学习器。这些学习器过于复杂，因而其风险高于正确权衡偏差与方差的模型。Hastie 等（2001）更详细地讨论了这一现象。

交叉验证能够更准确地估计风险，而不偏向更复杂的模型。我们的表述采用划分向量 $B_n\in\{0,1\}^n$，它将数据分为两组：训练集
$(O_i:B_n(i)=0)$，经验分布为 $P^0_{n,B_n}$；以及验证集
$(O_i:B_n(i)=1)$，经验分布为 $P^1_{n,B_n}$。
关于 $B_n$ 的分布取平均，得到交叉验证风险估计
\[
 \hat\theta_{\mathrm{CV}}(\bar Q)
 =E_{B_n}\theta(\bar Q(P^0_{n,B_n}),P^1_{n,B_n}).
\]
它给出交叉验证选择器
\[
 \hat k_n=\argmin_{k\in\{1,\ldots,K_n\}}\theta_{\mathrm{CV}}(\bar Q),
\]
即选择使交叉验证风险估计最小的学习器。由于对每一个划分向量 $B_n$，交叉验证都用不同的数据进行训练和风险估计，因此它具有重要的预言性质。

在适当条件下，就风险差比值（risk difference ratio）而言，交叉验证选择器在渐近意义上与预言选择器表现一样好：
\begin{equation}
 \frac{\theta_0(Q_{\hat k})-\theta_0(\bar Q_0)}
 {\theta_0(Q_{\tilde k})-\theta_0(\bar Q_0)}
 \xrightarrow{P}1.
\tag{28.1}
\end{equation}
也就是说，分子为交叉验证选择器的风险与最优风险之差，分母为预言选择器的风险与最优风险之差，二者之比依概率趋于 1。该结果的条件与证明见 Dudoit 和 van der Laan（2005）、van der Laan 和 Dudoit（2003）以及 van der Vaart 等（2006）。正是凭借这一性质，离散超级学习器在渐近意义上与其候选学习器中最好的一个表现一样好。下面将说明，这一性质在非参数自助样本上为何失效。
